\documentclass[aspectratio=169]{beamer}
\usepackage[T2A]{fontenc}
\usepackage[utf8]{inputenc}
\usepackage[russian]{babel}
\usepackage{helvet}
\renewcommand{\familydefault}{\sfdefault}
\usepackage{tikz}

% --- Настройка цветов (Linear Dark Style) ---
\definecolor{BgDark}{HTML}{121212} % Глубокий темный
\definecolor{BgLighter}{HTML}{1E1E1E} % Чуть светлее для акцентов
\definecolor{LinearCyan}{HTML}{00E5FF} % Драйвовый голубой акцент
\definecolor{TextWhite}{HTML}{F5F5F5} % Основной текст
\definecolor{TextGray}{HTML}{B0B0B0} % Второстепенный текст

% --- Применение цветов к Beamer ---
\setbeamercolor{background canvas}{bg=BgDark}
\setbeamercolor{normal text}{fg=TextWhite}
\setbeamercolor{frametitle}{fg=LinearCyan}
\setbeamercolor{itemize item}{fg=LinearCyan}
\setbeamercolor{itemize subitem}{fg=TextGray}

% --- Убираем лишнее ---
\setbeamertemplate{navigation symbols}{}
\setbeamertemplate{footline}{}

% --- Кастомный заголовок слайда (Linear Style) ---
\setbeamertemplate{frametitle}{
    \vspace{0.5cm}
    \begin{tikzpicture}[remember picture,overlay]
        \draw[LinearCyan, thick] (-0.5, -0.3) -- (2.5, -0.3);
        \draw[LinearCyan, thick] (-0.5, -0.3) -- (-0.5, 0.5);
    \end{tikzpicture}
    \vspace{-0.5cm}
    \Large\bfseries\insertframetitle\par
    \vspace{0.3cm}
}

\begin{document}

% --- ТИТУЛЬНЫЙ СЛАЙД ---
\begin{frame}
    \begin{tikzpicture}[remember picture,overlay]
        % Линейный графический акцент
        \draw[LinearCyan, line width=2pt] (current page.south west) ++(1,1) -- ++(0, 5) -- ++(4, 0);
        \draw[LinearCyan, line width=1pt, opacity=0.5] (current page.south west) ++(1.2,0.8) -- ++(0, 4.8) -- ++(3.8, 0);
    \end{tikzpicture}
    
    \vspace{2cm}
    \hspace{1.5cm}
    \begin{minipage}{0.8\textwidth}
        {\color{LinearCyan}\Huge \textbf{IndorTrafficPlan}}\\[0.5cm]
        {\color{TextWhite}\Large Проектирование ОДД на максималках}\\[0.3cm]
        {\color{TextGray}\small Цифровые решения для дорожного хозяйства от ИндорСофт}
    \end{minipage}
\end{frame}

% --- СЛАЙД 2: ПРОБЛЕМАТИКА И РЕШЕНИЕ ---
\begin{frame}{Хватит тратить часы на рутину}
    \Large
    Создание ПОДД не должно быть болью. 
    \vspace{0.5cm}
    
    \normalsize
    \textbf{IndorTrafficPlan} — это система, где ИИ и автоматизация работают за инженера. От исходного состояния до проектного решения — в единой цифровой среде.
    
    \vspace{0.5cm}
    \begin{itemize}
        \item Автоматическая расстановка ТСОДД.
        \item Мгновенное вычисление участков необеспеченной видимости.
        \item Сравнение «было/стало» в один клик.
    \end{itemize}
\end{frame}

% --- СЛАЙД 3: КИЛЛЕР-ФИЧИ ---
\begin{frame}{Драйв в каждом инструменте}
    \begin{columns}[T]
        \begin{column}{0.45\textwidth}
            {\color{LinearCyan}\large \textbf{01. Искусственный интеллект}}
            \vspace{0.2cm}\\
            \small Нейросети автоматически распознают стандартные дорожные знаки по видео. Меньше ручного ввода — больше точности.
            \vspace{0.5cm}\\
            {\color{LinearCyan}\large \textbf{02. Динамические чертежи}}
            \vspace{0.2cm}\\
            \small Разбивка на листы, спрямлённый план и топооснова. Настроил один раз — система сама перестроит чертёж при изменениях.
        \end{column}
        \begin{column}{0.45\textwidth}
            {\color{LinearCyan}\large \textbf{03. Индивидуальные знаки}}
            \vspace{0.2cm}\\
            \small Встроенный редактор \textit{IndorRoadSign}. Проектируйте знаки любой сложности по ГОСТам РФ и СНГ.
            \vspace{0.5cm}\\
            {\color{LinearCyan}\large \textbf{04. Генерация ведомостей}}
            \vspace{0.2cm}\\
            \small Пакет спецификаций по Приказу № 49 Минтранса РФ формируется одной командой.
        \end{column}
    \end{columns}
\end{frame}

% --- СЛАЙД 4: ИНТЕГРАЦИЯ ---
\begin{frame}{Бесшовная экосистема}
    \large
    Проект не живёт в вакууме. IndorTrafficPlan легко общается с внешним миром:
    
    \vspace{0.5cm}
    \normalsize
    \begin{itemize}
        \item \textbf{Вход:} Импорт из дорожных лабораторий (RDT Line, НПО Регион), GPS-треки, таблицы, данные из IndorCAD.
        \item \textbf{Подложки:} Интернет-карты, аэрофотосъёмка, DWG с координатной привязкой.
        \item \textbf{Выход:} Экспорт проектного решения напрямую в ГИС IndorRoad для дальнейшей эксплуатации.
    \end{itemize}
\end{frame}

% --- СЛАЙД 5: О КОМПАНИИ ---
\begin{frame}{Кто за этим стоит?}
    \begin{tikzpicture}[remember picture,overlay]
        \draw[LinearCyan, thick] (current page.north east) ++(-3,-1) -- ++(2,0) -- ++(0,-2);
    \end{tikzpicture}
    
    {\color{LinearCyan}\Large \textbf{ИндорСофт}} — 20+ лет в дорожной отрасли.
    \vspace{0.5cm}
    
    \begin{itemize}
        \item \textbf{Масштаб:} Более 300 000 км дорог в РФ и за рубежом проектируется и эксплуатируется на нашем ПО.
        \item \textbf{Импортозамещение:} 100\% в Реестре российского ПО. Поддержка Astra Linux и РЕД ОС.
        \item \textbf{Поддержка:} Мощное комьюнити, онлайн-курсы, форум и прямая связь с разработчиками.
    \end{itemize}
\end{frame}

% --- СЛАЙД 6: CALL TO ACTION ---
\begin{frame}
    \begin{center}
        \vspace{1.5cm}
        {\color{LinearCyan}\Huge \textbf{Готовы ускориться?}}\\[0.8cm]
        {\large Скачайте полнофункциональную версию на 30 дней бесплатно.}\\[1cm]
        
        \begin{tikzpicture}
            \node[draw=LinearCyan, thick, text=TextWhite, inner sep=0.4cm] (btn) {\textbf{indorsoft.ru/products/trafficplan}};
        \end{tikzpicture}\\[1cm]
        
        {\color{TextGray}\small 8 800 333-08-05 \quad | \quad info@indorsoft.ru}
    \end{center}
\end{frame}

\end{document}
